\section{Experimental Setup}

\subsection{Datasets and Protocols}

The primary evaluation uses the Tongji cross-session protocol in two directions:
\begin{itemize}
    \item S1->S2: training/gallery data from session 1 and probe data from session 2.
    \item S2->S1: training/gallery data from session 2 and probe data from session 1.
\end{itemize}
The main reported Tongji numbers are bidirectional averages over S1->S2 and S2->S1 with seeds 42, 2026, and 2705.

IITD within split is used as a secondary validation. The IITD result is treated as a scope and saturation check, not as the primary claim evidence.

\subsection{Compared Methods}

The paper compares the following methods:
\begin{itemize}
    \item B1: ResNet18 + cross entropy + supervised contrastive learning.
    \item B2: FixedGaborResNet18.
    \item B4: ResNet18 + ArcFace.
    \item B6: ResNet18 + BNNeck + ArcFace.
    \item B7: ResNet18 + BNNeck + ArcFace + light SupCon.
\end{itemize}

\subsection{Training Recipe}

The common training recipe used by the final compared runs is:
\begin{itemize}
    \item epochs: 60;
    \item learning rate: $1 \times 10^{-4}$;
    \item weight decay: $1 \times 10^{-4}$;
    \item embedding dimension: 256;
    \item Tongji seeds: 42, 2026, 2705.
\end{itemize}
B6 uses ArcFace with scale $s=30.0$ and margin $m=0.5$. B7 adds a light supervised contrastive term for ablation.

\subsection{Metrics}

The reported metrics are Rank-1, Rank-5, Macro-F1, EER, TAR@FAR=1e-2, TAR@FAR=1e-3, parameter count, FLOPs, and average inference latency. TAR@FAR=1e-3 is emphasized because it measures strict low-FAR verification behavior.

\subsection{Hardware}

The exact hardware and software environment should be filled in after a verified environment record is available. The current evidence files report latency values, but the current draft does not yet include a verified full hardware description.
